\begin{abstract}
Battle-royale survival unfolds under successive safe-zone contractions rather
than as a single post-match outcome. We reconstruct official PUBG telemetry as
a spatio-temporal database (STDB) aligned to indexed zone-change intervals of
moving teams, dynamic zones, spatial events, and team relations, and use it to
define a team-level next-phase elimination hazard task. From 21,824
graph-ready matches, we build
66-dimensional team-phase features and three-channel dynamic graphs encoding
proximity, combat, and shared safe-zone pressure. A feature-only gradient
boosting model already predicts the hazard well (AUC 0.765), but a
DynamicSTGNN over the dynamic graphs improves to AUC 0.778 and AP 0.681 on
134,245 aligned test cells. Representation ablations show that the
telemetry-based representation raises the score-summary baseline from AUC
0.649 to 0.765, while graph ablations and shuffle controls show that
dynamic team relations add smaller but significant signal beyond the
features. The result supports using phase-aligned spatio-temporal
representations for live game telemetry analysis.
\end{abstract}
